\section{A Father's Childhood Layoff and His Son's Earnings Profile}\label{sec:stability}

The parents' spells out of work are the one component of what is passed on whose share of the IGE rises with the son's
age, and the reason is that a father's layoff in his son's childhood costs the son more the older the son gets. At equal
parental income and equal values of every other component, a layoff in the son's teens is unrelated to his earnings at
25--29 and lowers them by 0.110$^{***}$ log points per layoff year at 45--49; a plant closure does neither. The cost is
in the son's wage, it grows within the same son, it follows layoffs during his childhood and not those after he is 20,
and it holds in every birth cohort and under one rule for recording job ends in every wave. Figure~\ref{fig:layoffbands}
shows it next to its two contrasts: plant closures, which carry no cost, and daughters, whose wages the same layoffs
raise (Section~\ref{sec:daughters}). This section establishes the profile; Section~\ref{sec:standsfor} asks what the
layoff stands for and Section~\ref{sec:career} where in the son's career the cost shows. Everything is descriptive: a
layoff is not assigned at random.

\begin{figure}[!htbp]
\centering
\includegraphics[width=0.92\textwidth]{paperB_layoff_bands.pdf}
\caption{A father's teen layoff or plant closure and his children's hourly wage, by the child's age (PSID)}
\label{fig:layoffbands}
\begin{minipage}{\textwidth}\footnotesize \textit{Notes:} Coefficient of the father's share of the child's ages 13--19
with a layoff or firing (sons and daughters) or a plant closure (sons) on the child's ln hourly wage in each age band,
at equal ln parental income and the father's teen-window ln wage, ln hours, unemployment, zero-hours and volatility
measures; single-age dummies; PSID weights; father present; 95\% confidence intervals clustered by 1968 family. One
layoff year is one seventh of a unit share. Code: \texttt{111.PaperBFigureLayoffBands.py}.\end{minipage}
\end{figure}

\subsection{The profile}


Table~\ref{tab:layoff_age} reports the partial coefficient behind the spells component: the coefficient of the
father's share of the son's ages 13--19 with a layoff or firing on the son's log earnings in each age band, holding log
parental income and every component of the split fixed (the parents' wages, hours, other spells and place; Appendix
Table~\ref{tab:passed_on}, panel A).\footnote{Father present; the mean layoff share is 0.044 and 12 percent of sons have
a layoff year: 383 father-present sons with a teen layoff year are observed at 25--29 and 135 at 45--49
(Table~\ref{tab:layoff_selection}). Layoff = the PSID's reason a job ended `laid off or fired', closure = `plant closed
or moved', in the interview year or the one before. Code: \texttt{88.LayoffCostByAge.py}
(\texttt{output/layoff\_cost\_by\_age.md}).} It is 0.146 ($p = 0.19$) at 25--29, falls in
every band, and reaches $-0.768^{***}$ ($p = 0.005$) at 45--49; the change between the two bands is $-0.914^{***}$
($p<0.001$). The share is over seven years, so one layoff year goes with earnings lower by 11 log points (0.110) at 45--49.
Five-year birth-cohort dummies leave the profile unchanged ($-0.739^{***}$ at 45--49, change $-0.866^{***}$), so it is
not that later-born sons are observed only young. Keeping the years with zero earnings, in a Poisson regression of
earnings in levels, the change is $-0.523^{**}$ (panel D of Table~\ref{tab:layoff_mechanism}), so the cost is not
confined to sons who work. The father's plant closures show no profile: 0.127 at 25--29, 0.127
at 45--49, a change of 0.000 (s.e.\ 0.258); the difference between the two changes, on common bootstrap draws, is
$-0.914^{***}$ ($p = 0.005$). Leaving out any one of the eleven five-year birth cohorts keeps the change between
$-0.82^{***}$ and $-0.99^{***}$. Nor is it one era or one way of asking. The PSID recorded why a job ended under
different rules: from 1968 to 1983 only if the father's job was less than a year old, from 1984 every job that ended
in the year, from 2003 a roster of jobs. With one rule applied to every wave---a job that ended within twelve months of
the interview, reported by the employed or the unemployed father---the change is $-1.392^{***}$ ($p = 0.003$) and the
within-son slope $-0.391^{**}$. The within-son slope is $-0.210^{*}$ for sons born before 1963, $-0.248$ ($p = 0.14$)
for those born 1963--75, who reach 45--49 in 2008--23, and $-0.547^{**}$ for those born 1976 and later; equality across
the three cohorts is not rejected ($p = 0.73$).\footnote{The one-rule share correlates 0.88 with the main one. By the
era of the father's teen-window layoffs the cost is clear for 1968--75 and from 1984, and imprecise for 1976--83, where
fewer than 30 treated sons reach 45--49. Code: \texttt{103.LayoffCohortRegime.py}.}

\begin{table}[!htbp]
\centering
\caption{A childhood layoff costs the son more as he ages (PSID)}
\label{tab:layoff_age}
\resizebox{\ifdim\width>\textwidth\textwidth\else\width\fi}{!}{%
\resizebox*{!}{\ifdim\totalheight>0.85\textheight 0.85\textheight\else\totalheight\fi}{%
\begin{tabular}{lccccccc}
\toprule
 & \begin{tabular}[b]{@{}c@{}}25--29\end{tabular} & \begin{tabular}[b]{@{}c@{}}30--34\end{tabular} & \begin{tabular}[b]{@{}c@{}}35--39\end{tabular} & \begin{tabular}[b]{@{}c@{}}40--44\end{tabular} & \begin{tabular}[b]{@{}c@{}}45--49\end{tabular} & \begin{tabular}[b]{@{}c@{}}50--55\end{tabular} & \begin{tabular}[b]{@{}c@{}}Change\\45--49 minus\\25--29\end{tabular} \\
 & (1) & (2) & (3) & (4) & (5) & (6) & (7) \\
\midrule
\multicolumn{8}{l}{\textit{A. Ln earnings, single-age dummies}} \\
Father laid off or fired, all sons & $0.146$ & $-0.006$ & $-0.110$ & $-0.241$ & $-0.768^{***}$ & $-0.543$ & $-0.914^{***}$ \\
 & (0.111) & (0.135) & (0.157) & (0.216) & (0.272) & (0.369) & (0.252) \\
Father's plant closed, all sons & $0.127$ & $0.143$ & $0.415^{**}$ & $0.229$ & $0.127$ & $-0.018$ & $0.000$ \\
 & (0.144) & (0.199) & (0.188) & (0.260) & (0.295) & (0.397) & (0.258) \\
Father laid off or fired, White sons & $0.158$ & $-0.033$ & $-0.097$ & $-0.267$ & $-0.818^{***}$ & $-0.621$ & $-0.975^{***}$ \\
 & (0.118) & (0.149) & (0.174) & (0.232) & (0.297) & (0.442) & (0.274) \\
Father laid off or fired, Black sons & $-0.218$ & $-0.064$ & $-0.379$ & $-0.107$ & $-0.351$ & $-0.133$ & $-0.133$ \\
 & (0.279) & (0.341) & (0.453) & (0.473) & (0.444) & (0.321) & (0.431) \\
\addlinespace
\multicolumn{8}{l}{\textit{B. Ln earnings, with five-year birth-cohort dummies}} \\
Father laid off or fired, all sons & $0.127$ & $-0.013$ & $-0.116$ & $-0.245$ & $-0.739^{***}$ & $-0.543$ & $-0.866^{***}$ \\
 & (0.108) & (0.140) & (0.159) & (0.218) & (0.260) & (0.351) & (0.242) \\
Father's plant closed, all sons & $0.106$ & $0.121$ & $0.412^{**}$ & $0.220$ & $0.086$ & $-0.038$ & $-0.020$ \\
 & (0.146) & (0.200) & (0.189) & (0.259) & (0.286) & (0.394) & (0.253) \\
Father laid off or fired, White sons & $0.152$ & $-0.037$ & $-0.097$ & $-0.273$ & $-0.774^{***}$ & $-0.614$ & $-0.926^{***}$ \\
 & (0.118) & (0.154) & (0.175) & (0.235) & (0.284) & (0.419) & (0.263) \\
Father laid off or fired, Black sons & $-0.334$ & $-0.037$ & $-0.383$ & $-0.096$ & $-0.408$ & $-0.094$ & $-0.075$ \\
 & (0.286) & (0.342) & (0.451) & (0.426) & (0.437) & (0.343) & (0.441) \\
\addlinespace
\multicolumn{8}{l}{\textit{C. Which outcome: father laid off or fired (columns 25--29, 45--49 and the change)}} \\
All sons: ln annual earnings & $0.146$ &  &  &  & $-0.768^{***}$ &  & $-0.914^{***}$ \\
 & (0.111) &  &  &  & (0.272) &  & (0.252) \\
All sons: ln hourly wage & $0.046$ &  &  &  & $-0.557^{***}$ &  & $-0.603^{***}$ \\
 & (0.089) &  &  &  & (0.213) &  & (0.203) \\
All sons: ln annual hours (employed) & $0.100^{*}$ &  &  &  & $-0.211$ &  & $-0.311^{**}$ \\
 & (0.059) &  &  &  & (0.143) &  & (0.139) \\
All sons: employed & $0.050^{*}$ &  &  &  & $0.083$ &  & $0.033$ \\
 & (0.026) &  &  &  & (0.078) &  & (0.080) \\
All sons: his share of years unemployed & $0.008$ &  &  &  & $0.030$ &  & $0.022$ \\
 & (0.075) &  &  &  & (0.086) &  & (0.103) \\
\bottomrule
\end{tabular}}}
\par\smallskip
\begin{minipage}{\textwidth}\footnotesize \textit{Notes:} Coefficient of the father's share of the son's ages 13--19 with a layoff or firing (PSID 'laid off / fired') or a plant closure ('plant closed / moved') on the son's outcome in each age band, holding ln parental income and every component of Table~\ref{tab:passed_on}, panel A, fixed; single-age dummies; PSID weights. The share is over seven years, so one year with a layoff is one seventh of the coefficient. Father present; mean layoff share 0.044, 12 percent of sons have a layoff year. Standard errors from 200 cluster-bootstrap draws of 1968 families. Source: \texttt{88.LayoffCostByAge.py}; table built by \texttt{90.PaperBRaceNeutralTables.py}. $^{***}$ $p<0.01$, $^{**}$ $p<0.05$, $^{*}$ $p<0.1$.\end{minipage}
\end{table}


The cost is in the wage. In panel C of Table~\ref{tab:layoff_age} the coefficient on the son's log hourly wage changes by $-0.603^{***}$ ($p = 0.003$)
between 25--29 and 45--49, two thirds of the change in earnings, and that on his annual hours when employed by
$-0.311^{**}$ ($p = 0.025$); his employment ($0.033$) and his share of years unemployed ($0.022$) do not change. A son
whose father was laid off in his teens is as often employed as other sons at 45--49, works fewer hours when employed and
is paid less per hour.

\paragraph{Within the same son.} The profile across bands compares different sons at different ages. Table~\ref{tab:layoff_fe}
follows the same son: with son fixed effects, the father's layoff share interacted with the son's age measures how his
earnings profile bends. A unit share lowers the growth of the same son's log earnings by 0.236$^{**}$ ($p = 0.022$) per
ten years of age, 0.034 a decade for one layoff year, with every component, parental income, the family-structure flags
and cohort dummies also interacted with age. His hourly wage bends by $-0.109$ ($p = 0.18$), so within the same son
the bend is clearer in earnings than in the wage, and it sits in his self-employed years
(Section~\ref{sec:career}). Closures do not bend it ($-0.010$, s.e.\ 0.119). Because each son is
compared with himself, a fixed difference in the level of earnings between sons whose fathers were and were not laid
off cannot produce the result. A fixed difference whose return grows with age---in skill, say---can, and
Section~\ref{sec:standsfor} asks what the layoff stands for.

\paragraph{Timing.} Layoffs after the son is 20 carry no cost; whether teen layoffs cost more than those at 6--12 the data
cannot say. Adding the father's layoff shares at the son's ages 6--12 and 20--26 leaves the teen coefficient at
$-0.322^{***}$ ($p = 0.010$), against $-0.081$ (s.e.\ 0.109) at 6--12 and 0.037 (s.e.\ 0.102) at 20--26 (columns 2, 4
and 6). The teen coefficient differs from the one at 20--26 ($-0.359^{**}$, $p = 0.043$) but not from the one at 6--12
($-0.241$, $p = 0.18$), and equality of all three is rejected only at 13 percent (panel B of
Table~\ref{tab:layoff_selection}).\footnote{These windows need the father as a head at the son's 6--12 and 20--26,
which keeps 2,171 of the 4,644 sons; on them the teen coefficient alone is the same as on all sons.}

\begin{table}[!htbp]
\centering
\caption{Within the same son: the earnings profile and the father's childhood layoffs (PSID)}
\label{tab:layoff_fe}
\resizebox{\ifdim\width>\textwidth\textwidth\else\width\fi}{!}{%
\resizebox*{!}{\ifdim\totalheight>0.85\textheight 0.85\textheight\else\totalheight\fi}{%
\begin{tabular}{lcccccc}
\toprule
 & \begin{tabular}[b]{@{}c@{}}White\end{tabular} & \begin{tabular}[b]{@{}c@{}}White,\\with timing\end{tabular} & \begin{tabular}[b]{@{}c@{}}Black\end{tabular} & \begin{tabular}[b]{@{}c@{}}Black,\\with timing\end{tabular} & \begin{tabular}[b]{@{}c@{}}All\end{tabular} & \begin{tabular}[b]{@{}c@{}}All,\\with timing\end{tabular} \\
 & (1) & (2) & (3) & (4) & (5) & (6) \\
\midrule
Layoff share, ages 13--19, $\times$ age & $-0.229^{**}$ & $-0.338^{**}$ & $-0.306^{**}$ & $-0.068$ & $-0.236^{**}$ & $-0.322^{***}$ \\
 & (0.113) & (0.135) & (0.126) & (0.186) & (0.103) & (0.124) \\
Closure share, ages 13--19, $\times$ age & $-0.007$ & $0.042$ & $0.009$ & $0.395^{*}$ & $-0.010$ & $0.045$ \\
 & (0.128) & (0.119) & (0.191) & (0.219) & (0.119) & (0.116) \\
Layoff share, ages 6--12, $\times$ age &  & $-0.097$ &  & $-0.021$ &  & $-0.081$ \\
 &  & (0.117) &  & (0.226) &  & (0.109) \\
Layoff share, ages 20--26, $\times$ age &  & $0.061$ &  & $0.073$ &  & $0.037$ \\
 &  & (0.105) &  & (0.226) &  & (0.102) \\
Closure share, ages 6--12, $\times$ age &  & $0.101$ &  & $-0.307$ &  & $0.054$ \\
 &  & (0.132) &  & (0.216) &  & (0.125) \\
Closure share, ages 20--26, $\times$ age &  & $-0.017$ &  & $-0.369^{*}$ &  & $-0.045$ \\
 &  & (0.177) &  & (0.223) &  & (0.167) \\
Ln parental income $\times$ age & $0.060$ & $0.054$ & $-0.011$ & $-0.069$ & $0.033$ & $0.051$ \\
 & (0.041) & (0.057) & (0.048) & (0.099) & (0.033) & (0.054) \\
\bottomrule
\end{tabular}}}
\par\smallskip
\begin{minipage}{\textwidth}\footnotesize \textit{Notes:} Son-years 25--55, ln annual earnings (valid years) with son fixed effects; single-age dummies; every component of Table~\ref{tab:passed_on}, panel A, ln parental income, the family-structure flags and five-year cohort dummies, each interacted with (age $-$ 40)/10. A coefficient is the change per ten years of age in the same son's ln earnings per unit share (one layoff year = one seventh). 'With timing' adds the father's layoff and closure shares at the son's ages 6--12 and 20--26 (his own head-years in the PSID) $\times$ age. PSID weights; standard errors clustered by 1968 family. Source: \texttt{IGE/88.LayoffCostByAge.py}; table built by \texttt{IGE/90.PaperBRaceNeutralTables.py}. $^{***}$ $p<0.01$, $^{**}$ $p<0.05$, $^{*}$ $p<0.1$.\end{minipage}
\end{table}


The closure is the contrast that disciplines the reading of this profile. A plant closure displaces workers whatever their quality,
so at equal parental income it says little about the father, and the income it costs is already in parental income
\citep{gibbons1991layoffs}; a layoff selects workers. That closures carry no profile, and at equal parental income even
raise the son's wage (Table~\ref{tab:four_by_four}), is the pattern expected if the layoff coefficient carries what
layoffs bring with them beyond lost income: the father's subsequent instability, the family's response or what the
layoff reveals about the father's job. The design cannot separate these, and it does not say what a randomly assigned
layoff would do. The closure null is also imprecise: its 95 percent interval admits a closure cost up to 59 percent of
the layoff's, so what it rules out is a closure cost as large as the layoff's.

Self-employed fathers do not make the profile. A father who works for himself cannot report a layoff, so his son sits
among the sons of fathers never laid off, and a business that fails is recorded under the same code as a plant
closure; and sons of self-employed fathers are more often in their father's line of work after 35
(Appendix~\ref{sec:fathers}), so their earnings could rise late and lend the layoff a late cost. Without the sons of
fathers self-employed in at least half of the son's teen waves---18.5 percent of father-present sons---the change
between 25--29 and 45--49 is $-0.880^{***}$ ($p<0.001$) and the within-son slope $-0.295^{***}$ ($p = 0.001$), against
$-0.914^{***}$ and $-0.236^{**}$ on all sons; without fathers self-employed in any teen wave they are $-0.912^{***}$ and
$-0.289^{***}$ ($p = 0.005$), and the layoff's change exceeds the closure's by $-1.102^{***}$ ($p = 0.001$). Holding the
father's self-employment and its interaction with age fixed leaves them at $-0.896^{***}$ and
$-0.233^{**}$.\footnote{Self-employed = the head works for himself, alone or also for others, at the interview, in the
father's head-waves at the son's 13--19 (asked in every wave from 1968). Self-employed fathers have a mean layoff share of
0.020 against 0.051; 34 of the 430 father-present sons with a teen layoff year have a self-employed father. The PSID's
closure code bundles plants that closed or moved with employers that went out of business. The sons--daughters
difference in the wage change is 0.773$^{**}$ ($p = 0.010$) without self-employed fathers (Section~\ref{sec:daughters}).
Code: \texttt{119.FatherSelfEmployment.py} (\texttt{output/father\_self\_employment.md}).}

The age profile reconciles the result with the evidence that a father's displacement costs his children little
\citep{oreopoulos2008intergenerational, hilger2016parental}. Those designs score children's earnings before their
mid-thirties; at those ages this paper too finds no cost of a father's layoff, and plant closures, the displacement
those designs isolate best, carry none at any age. The cost appears after 40, beyond the ages they observe. Three further
checks hold: in ranks, the son falls 22.2$^{***}$ percentiles between 25--29 and 45--49 per unit share, about three per
layoff year; the minimum detectable change at 45--49 is 0.68, against an estimate of 0.91; and the six tests of the cost
reported in this section survive family-wise corrections, with adjusted $p$-values from 0.001 to 0.044 (Holm) and from
0.003 to 0.035 (Romano--Wolf step-down on a joint cluster bootstrap) (panels B and
D of Table~\ref{tab:layoff_referee}).\footnote{A father figure's displacement of any kind at the son's ages 0--17, compared with sons never
exposed whose families had the same income before the event, is associated with earnings lower by 0.079$^{***}$ at
25--34 and 0.053 (s.e.\ 0.043) at 35--44 (Appendix Table~\ref{tab:fatherdispage}). That design pools closures with layoffs,
spreads exposure over childhood and stops at 44, before the layoff cost appears; its flat profile is consistent with the
results here. Code: \texttt{49.PaperBPolicyShocks.py}.}

\paragraph{Timing as the control group.} Among sons whose father was laid off at some time, those whose father was laid
off during their childhood fall behind those whose father was laid off only when they were 20--26. Restricting the
within-son design to the 547 sons whose father had a layoff at 6--12, 13--19 or 20--26 compares fathers who all had a
layoff and differ in when it fell. Relative to sons whose father's layoffs came only at 20--26, the earnings of sons
with a teen layoff grow more slowly by 0.157$^{**}$ ($p = 0.021$) a decade, and those of sons with a layoff only at
6--12 by 0.110$^{*}$ ($p = 0.081$); the two childhood groups do not differ ($-0.047$, $p = 0.45$), and the teen group's
profile falls behind by 0.329$^{**}$ at 45--49 (panel C of Table~\ref{tab:layoff_selection}). A trait common to fathers
who are ever laid off cannot produce a cost that depends on when the layoff fell; a trait of fathers laid off early in
the son's life, or of younger fathers, can.\footnote{Fathers laid off only at the son's 20--26 are older at the layoff.
The comparison removes what all laid-off fathers share, not what distinguishes early from late layoffs.}

\paragraph{Mothers who head their families.} The layoff analysis follows fathers, so it leaves out the families headed
by a mother, the main recipients of the safety net. The PSID asks female heads the same job-end question, and a
mother's layoffs during her children's teens go with a delayed cost, for her daughters. Among the 2,843 children of families
headed by a woman in at least three of their teen waves (74 percent Black; 507 with a teen layoff year of their
mother), a unit of her teen layoff share bends the child's earnings profile by $-0.368^{**}$ ($p = 0.029$) a decade and
the wage by $-0.236^{**}$ ($p = 0.030$) within the same child, about the size of the father's cost for sons. The comparison
across bands is imprecise ($-0.215$ between 25--29 and 40--44, s.e.\ 0.275), and by sex the within-child cost is
$-0.079$ ($p = 0.64$) for sons and $-0.561^{**}$ ($p = 0.027$) for daughters: in these families the pooled cost is the
daughters'. It is not a pattern by the parent's sex: in two-parent families, where wives report their own job ends
from 1984, a mother's teen layoffs go with nothing for her sons ($-0.049$, $p = 0.64$, wage) and a higher wage for her
daughters ($+0.228^{*}$, $p = 0.056$), as the father's layoffs do (\texttt{139.MothersAndAFDC.py}).\footnote{Female heads' job-end reports with
the fathers' definitions and dating rules; layoff share over the child's teen waves with a female head (three or
more), at equal log parental income and closure share; sons and daughters as heads or spouses at 25--55; 190 treated
children are observed at 40--44. Code: \texttt{131.MotherLayoffs.py} (\texttt{output/mother\_layoffs.md}).}

\begin{table}[!htbp]
\centering
\caption{Is the delayed cost of a father's layoff selection? Timing, comparison groups and splits (PSID)}
\label{tab:layoff_selection}
\resizebox{\ifdim\width>\textwidth\textwidth\else\width\fi}{!}{%
\resizebox*{!}{\ifdim\totalheight>0.85\textheight 0.85\textheight\else\totalheight\fi}{%
\begin{tabular}{lccc}
\toprule
 & \begin{tabular}[b]{@{}c@{}}Pooled\end{tabular} & \begin{tabular}[b]{@{}c@{}}White\end{tabular} & \begin{tabular}[b]{@{}c@{}}Black\end{tabular} \\
 & (1) & (2) & (3) \\
\midrule
\multicolumn{4}{l}{\textit{A. Counts and the closure contrast}} \\
Treated sons at 25--29 / 45--49 (father present) & 383 / 135 & 265 / 98 & 118 / 37 \\
Layoff minus closure change, 45--49 minus 25--29 & $-0.914^{***}$ & $-0.962^{***}$ &  \\
 & (0.326) & (0.338) &  \\
\addlinespace
\multicolumn{4}{l}{\textit{B. Three windows of the son's age, within son}} \\
Layoff share, ages 6--12, $\times$ age & $-0.081$ & $-0.097$ & $-0.021$ \\
 & (0.109) & (0.117) & (0.226) \\
Layoff share, ages 13--19, $\times$ age & $-0.322^{***}$ & $-0.338^{**}$ & $-0.068$ \\
 & (0.124) & (0.135) & (0.186) \\
Layoff share, ages 20--26, $\times$ age & $0.037$ & $0.061$ & $0.073$ \\
 & (0.102) & (0.105) & (0.226) \\
13--19 minus 6--12 & $-0.241$ & $-0.241$ & $-0.047$ \\
 & (0.180) & (0.196) & (0.335) \\
13--19 minus 20--26 & $-0.359^{**}$ & $-0.399^{**}$ & $-0.141$ \\
 & (0.177) & (0.188) & (0.341) \\
All three equal: $\chi^2(2)$ [$p$] & 4.12 [0.128] & 4.52 [0.104] & 0.18 [0.915] \\
Sons & 2,171 & 1,603 & 568 \\
\addlinespace
\multicolumn{4}{l}{\textit{C. Timing as the control: sons whose father was laid off at some time (reference: only at 20--26)}} \\
Layoff at 13--19 (T) $\times$ age & $-0.157^{**}$ & $-0.167^{**}$ & $0.062$ \\
 & (0.068) & (0.074) & (0.113) \\
Layoff at 6--12 only (E) $\times$ age & $-0.110^{*}$ & $-0.122^{*}$ & $0.072$ \\
 & (0.063) & (0.070) & (0.133) \\
T minus E & $-0.047$ & $-0.045$ & $-0.010$ \\
 & (0.062) & (0.067) & (0.143) \\
Sons T / E / L & 256 / 164 / 127 & 195 / 113 / 98 & 61 / 51 / 29 \\
\bottomrule
\end{tabular}}}
\par\smallskip
\begin{minipage}{\textwidth}\footnotesize \textit{Continued on the next table; notes below its last part.}\end{minipage}
\end{table}
\begin{table}[!htbp]
\centering
\addtocounter{table}{-1}
\caption{Is the delayed cost of a father's layoff selection? Timing, comparison groups and splits (PSID) (continued)}
\resizebox{\ifdim\width>\textwidth\textwidth\else\width\fi}{!}{%
\resizebox*{!}{\ifdim\totalheight>0.85\textheight 0.85\textheight\else\totalheight\fi}{%
\begin{tabular}{lccc}
\toprule
 & \begin{tabular}[b]{@{}c@{}}Pooled\end{tabular} & \begin{tabular}[b]{@{}c@{}}White\end{tabular} & \begin{tabular}[b]{@{}c@{}}Black\end{tabular} \\
 & (1) & (2) & (3) \\
\midrule
\multicolumn{4}{l}{\textit{D. What the layoff is and what it does: the teen layoff share split into parts (pooled)}} \\
Layoff share, small fall in family income, $\times$ age & $-0.473^{***}$ &  &  \\
 & (0.144) &  &  \\
Layoff share, large fall in family income, $\times$ age & $-0.096$ &  &  \\
 & (0.132) &  &  \\
\quad equal: $p$ & 0.079 &  &  \\
Layoff share, short unemployment (< 500 hours), $\times$ age & $-0.338^{***}$ &  &  \\
 & (0.119) &  &  \\
Layoff share, long unemployment (500+ hours), $\times$ age & $0.225$ &  &  \\
 & (0.265) &  &  \\
\quad equal: $p$ & 0.064 &  &  \\
Layoff share, one teen layoff year, $\times$ age & $0.295$ &  &  \\
 & (0.378) &  &  \\
Layoff share, two or more, $\times$ age & $-0.268^{**}$ &  &  \\
 & (0.106) &  &  \\
\quad equal: $p$ & 0.145 &  &  \\
Layoff share, no UI received, $\times$ age & $-0.211$ &  &  \\
 & (0.255) &  &  \\
Layoff share, UI received, $\times$ age & $-1.014^{**}$ &  &  \\
 & (0.484) &  &  \\
\quad equal: $p$ & 0.095 &  &  \\
Layoff share $\times$ age, + ln parental income at 20--26 & $-0.232^{**}$ &  &  \\
 & (0.102) &  &  \\
\bottomrule
\end{tabular}}}
\par\smallskip
\begin{minipage}{\textwidth}\footnotesize \textit{Notes:} Son-years 25--55 with a valid wage; ln earnings. Within son: son fixed effects, single-age dummies, ln parental income, every component of Table~\ref{tab:passed_on} panel A, the family-structure flags and five-year cohort dummies, each $\times$ (age $-$ 40)/10; coefficients per unit share per ten years of age; standard errors clustered by 1968 family. Panel A, second row: band coefficients at equal parental income and components, common bootstrap draws (200). Panel B: sons whose father is observed as a head at their 6--12 and 20--26. Panel C: T = a layoff at 13--19, E = at 6--12 and not 13--19, L = only at 20--26. Panel D: each split replaces the layoff share by parts that add up to it (a third part, not measured, is in the regression and not shown); the fall in family income is the layoff wave against the son's layoff-free teen waves, split at the median (-0.084); hours unemployed are measured to 1993 and UI receipt in 28\% of layoff years. Source: \texttt{91.LayoffSelection.py}; table built by \texttt{100.PaperBFableTables.py}. $^{***}$ $p<0.01$, $^{**}$ $p<0.05$, $^{*}$ $p<0.1$.\end{minipage}
\end{table}


\subsection{By race}

The rise of the spells component with age is White. On common bootstrap draws, its contribution to the earnings IGE
rises between 25--29 and 45--49 by 0.066$^{**}$ ($p = 0.010$) for White sons and changes by $-0.042$ ($p = 0.25$) for
Black sons, a difference of 0.107$^{**}$ ($p = 0.018$); with birth-cohort dummies 0.120$^{***}$ ($p = 0.006$). The
difference is in the father's layoffs (0.040$^{**}$, $p = 0.028$) and the mother's unemployment (0.065$^{**}$,
$p = 0.049$). It is not shown that a layoff costs Black sons less. Within the same son a layoff bends the Black
earnings profile by $-0.306^{**}$ ($p = 0.015$) per ten years against $-0.229^{**}$ for White sons
(Table~\ref{tab:layoff_fe}), but the Black estimate rests on 37 treated sons at 45--49 and falls to $-0.068$ with the
timing controls and to $+0.062$ in the timing-as-control design: the cost is not distinguishable by race.
It is that Black fathers' layoffs are not concentrated in low-income families: the gradient of the layoff share in log
parental income is $-0.046^{***}$ for White fathers at 45--49 and $-0.010^{*}$ for Black fathers, so a cost that rises
with age raises the White IGE and leaves the Black IGE unchanged.\footnote{Code: \texttt{81.StabilityOpeningByRace.py}
(\texttt{output/stability\_opening\_by\_race.md}).} The racial difference in what is passed on through spells is in who
is laid off, not in what a layoff does.


